\documentclass[../main.tex]{subfiles}
\begin{document}

% --------------------------------------------------
\section{フォルダ構成}
% --------------------------------------------------

このテンプレートは, 本文, 説明文, スタイルファイル, 文献データ, 画像ファイルを役割ごとに分けている. 
最初はファイル数が多く見えるかもしれないが, 基本的には\texttt{main.tex}を入口として, 必要に応じて\texttt{doc/}や\texttt{sty/}の中を編集すればよい. 

\subsection{全体像}

主なフォルダ構成は次の通りである. 

\begin{verbatim}
note/
├── main.tex
├── main.synctex.gz
├── main.pdf
├── sty/
│     ├── note.sty
│     ├── core.sty
:      :
│     └── bib.sty
├── doc/
│     ├── structure.tex
│     ├── note.tex
:      :
│     └── bib.tex
├── bib/
│     └── refs.bib
├── img/
│      └── coordinate_transform.tex
├── ref
└── aux
\end{verbatim}

\subsection{入口となるファイル}

\texttt{main.tex}は, 文書全体を組み立てるための入口である. 
文書クラス, \texttt{note.sty}の読み込み, 数式作用素の定義, 文献ファイルの登録, タイトル情報, 本文ファイルの読み込みをここに書いている. 

本文を増やしたい場合は, \texttt{doc/}の中に新しい\texttt{tex}ファイルを作り, \texttt{main.tex}に
\begin{center}
	\verb|\subfile{doc/<ファイル名>.tex}|
\end{center}
を追加する. 
既存の説明文もこの方法で分割している. 

\subsection{スタイルファイル}

\texttt{sty/}には, 見た目やマクロを定義するスタイルファイルを置いている. 
利用者が通常読み込むのは\texttt{sty/note.sty}だけでよい. 
\texttt{note.sty}が他のスタイルファイルを順に読み込むため, 個別に\texttt{core.sty}や\texttt{math.sty}を読み込む必要はない. 

各スタイルファイルの役割は, おおまかに次の通りである. 
\begin{itemize}
	\item \texttt{core.sty} : 基本パッケージ, 色, 数式表示の共通設定
	\item \texttt{layout.sty} : 背景色, ページスタイル, 目次, 見出し, 箇条書き
	\item \texttt{math.sty} : 数式用マクロ, 括弧, アクセント, 特殊記号
	\item \texttt{theorem.sty} : 定理系環境, 証明環境, Pythonコード環境
	\item \texttt{code.sty} : コード掲載機能を拡張するための予約ファイル
	\item \texttt{ref.sty} : \texttt{zref-clever}による日本語参照設定
	\item \texttt{bib.sty} : \texttt{biblatex}による文献表示設定
\end{itemize}

\subsection{説明文ファイル}

\texttt{doc/}には, このテンプレートの説明文を置いている. 
各\texttt{tex}ファイルは, 対応するスタイルファイルの使い方を説明するためのものである. 
たとえば\texttt{doc/math.tex}は\texttt{sty/math.sty}の説明, \texttt{doc/bib.tex}は\texttt{sty/bib.sty}の説明になっている. 

各ファイルは\texttt{subfiles}を使っているため, 親ファイルである\texttt{main.tex}から読み込むことも, 必要に応じて単独でコンパイルすることもできる. 
新しい説明ファイルを作る場合は, 次の形を雛形にするとよい. 
\begin{verbatim}
\documentclass[../main.tex]{subfiles}
\begin{document}

% --------------------------------------------------
\section{見出し}
% --------------------------------------------------

本文を書く. 

\end{document}
\end{verbatim}

\subsection{文献と画像}

\texttt{bib/refs.bib}には, \texttt{biblatex}で使う文献データを置く. 
文献を追加したい場合はこのファイルにBibTeX形式で項目を追加し, 本文中では\verb|\cite{<引用ラベル>}|で引用する. 

\texttt{img/}には, 本文中で使う図や画像を置く. 
現在は\texttt{coordinate\_transform.tex}のように, TikZで描いた図を\verb|\input|で読み込む形にしている. 
例えば\texttt{doc/note.tex}では, \texttt{img/}フォルダにある\texttt{coordinate\_transform.tex}を次のように読み込んでいる. 
\begin{verbatim}
\begin{figure}[tbp]
	\centering
	\resizebox{\linewidth}{!}{%
		\usetikzlibrary{positioning,calc}
\begin{tikzpicture}
	% \draw[thin, dotted] (0, 0) grid(23, 11); % 補助grid

	% 位相空間 X
	\draw[tension=0.7,thin] plot[smooth cycle] coordinates{
		(0.5,6) (2.5,2) (6,0.5) (8.5,2.5) (10.5,6) (9.5,9) (6,10.5) (3,8)
	}
	(2,9) node{$X$};
	
	% U
	\begin{scope}[local bounding box=U]
		\draw[densely dashed,line width=1pt,color=ORANGE] (6, 7) circle[radius=2.5] (8,8) node{$U$};
	\end{scope}

	% Vの中の座標系
	\begin{scope}
		\clip (6,7) circle[radius=2.5];
		\draw[->,thick,color=ORANGE,rotate around={-20:(6,7)}] (2,6) --(10,6) node[right]{};
		\draw[->,thick,color=ORANGE,rotate around={-20:(6,7)}] (5,4) --(5,11) node[right]{};
		\draw[ORANGE!40, very thin, rotate around={-20:(6,7)},step=1.67,shift={(0,-0.67)}] (2,4) grid(10,11);
	\end{scope}

	% V
	\begin{scope}[local bounding box=V]
		\draw[densely dashed,line width=1pt,color=GREEN] (6, 4) circle[radius=2.5] (8,3) node{$V$};
	\end{scope}

	% Uの中の座標系
	\begin{scope}
		\clip (6,4) circle[radius=2.5];
		\draw[->,thick,color=GREEN,rotate around={-20:(6,4)}] (2,3) --(10,3) node[right]{};
		\draw[->,thick,color=GREEN,rotate around={-20:(6,4)}] (5,1) --(5,8) node[right]{};
		\draw[GREEN!40, very thin, rotate around={-20:(6,4)},step=1.67,shift={(0,-0.33)}] (2,1) grid(10,8);
	\end{scope}

	% 点p
	\coordinate (p) at (6, 5.5);
	\draw[color=BLUE,fill=BLUE] (p) circle[radius=0.1] node[anchor=south]{$p$};

	% R^mの座標系
	\draw[ORANGE!40, very thin, step=1] (15.1,6.5) grid(21.9,10.5);
	\draw[->,thick,color=ORANGE] (15,8) --(22,8) node[right]{$x_1$};
	\draw[->,thick,color=ORANGE] (18,6.5) --(18,10.5) node[right]{$x_2$};

	% Udash
	\begin{scope}[local bounding box=Udash]
		\draw[densely dashed,line width=1pt,color=ORANGE] (18.5, 8.5) circle[radius=1.5] (20, 10) node{$\varphi(U)$};
	\end{scope}

	% 2つの円の共通部分
	\begin{scope}
		\clip (18.5,8.5) circle[radius=1.5];
		\draw[densely dashed,line width=1pt,color=GREEN] (19.3,6.9) circle[radius=1.5];
	\end{scope}

	% 点phi(p)
	\coordinate (p) at (18.8, 7.7);
	\draw[color=BLUE,fill=BLUE] (p) circle[radius=0.1] node[anchor=west]{$\varphi(p)$};

	% R^mの座標系
	\draw[GREEN!40, very thin, step=1] (15.1,0.5) grid(21.9,4.5);
	\draw[->,thick,color=GREEN] (15,2) --(22,2) node[right]{$y_1$};
	\draw[->,thick,color=GREEN] (18,0.5) --(18,4.5) node[right]{$y_2$};

	% Vdash
	\begin{scope}[local bounding box=Vdash]
		\draw[densely dashed,line width=1pt,color=GREEN] (18.5, 2.5) circle[radius=1.5] (20, 2) node{$\psi(V)$};
	\end{scope}

	% 2つの円の共通部分
	\begin{scope}
		\clip (18.5,2.5) circle[radius=1.5];
		\draw[densely dashed,line width=1pt,color=ORANGE] (17.9,4.3) circle[radius=1.5];
	\end{scope}

	% 点psi(p)
	\coordinate (p) at (18.3, 3.4);
	\draw[color=BLUE,fill=BLUE] (p) circle[radius=0.1] node[anchor=east]{$\psi(p)$};

	% 繋ぎ線
	\draw[->,thick,shorten <=10pt, shorten >=10pt] (Udash.south) --(Vdash.north) node[midway, right]{$\psi\circ\varphi^{-1}$};
	\draw[->,thick,shorten <=10pt, shorten >=10pt] (Udash.west) --(U.east) node[midway, above]{$\varphi^{-1}$};
	\draw[->,thick,shorten <=10pt, shorten >=10pt] (Vdash.west) --(V.east) node[midway, below]{$\psi^{-1}$};

\end{tikzpicture}

	}
	\caption{座標変換の模式図}
	\label{fig:sample-coordinate-transform}
\end{figure}
\end{verbatim}
このように, 図形を描くための\texttt{tex}ファイルを\texttt{img/}に置いておくと, 本文側では\verb|\input|で呼び出すだけでよい. 
画像ファイルを使う場合も, 同じフォルダに置いて\verb|\includegraphics|で読み込めばよい. 

\begin{figure}[htbp]
	\centering
	\resizebox{0.8\linewidth}{!}{%
		\usetikzlibrary{positioning,calc}
\begin{tikzpicture}
	% \draw[thin, dotted] (0, 0) grid(23, 11); % 補助grid

	% 位相空間 X
	\draw[tension=0.7,thin] plot[smooth cycle] coordinates{
		(0.5,6) (2.5,2) (6,0.5) (8.5,2.5) (10.5,6) (9.5,9) (6,10.5) (3,8)
	}
	(2,9) node{$X$};
	
	% U
	\begin{scope}[local bounding box=U]
		\draw[densely dashed,line width=1pt,color=ORANGE] (6, 7) circle[radius=2.5] (8,8) node{$U$};
	\end{scope}

	% Vの中の座標系
	\begin{scope}
		\clip (6,7) circle[radius=2.5];
		\draw[->,thick,color=ORANGE,rotate around={-20:(6,7)}] (2,6) --(10,6) node[right]{};
		\draw[->,thick,color=ORANGE,rotate around={-20:(6,7)}] (5,4) --(5,11) node[right]{};
		\draw[ORANGE!40, very thin, rotate around={-20:(6,7)},step=1.67,shift={(0,-0.67)}] (2,4) grid(10,11);
	\end{scope}

	% V
	\begin{scope}[local bounding box=V]
		\draw[densely dashed,line width=1pt,color=GREEN] (6, 4) circle[radius=2.5] (8,3) node{$V$};
	\end{scope}

	% Uの中の座標系
	\begin{scope}
		\clip (6,4) circle[radius=2.5];
		\draw[->,thick,color=GREEN,rotate around={-20:(6,4)}] (2,3) --(10,3) node[right]{};
		\draw[->,thick,color=GREEN,rotate around={-20:(6,4)}] (5,1) --(5,8) node[right]{};
		\draw[GREEN!40, very thin, rotate around={-20:(6,4)},step=1.67,shift={(0,-0.33)}] (2,1) grid(10,8);
	\end{scope}

	% 点p
	\coordinate (p) at (6, 5.5);
	\draw[color=BLUE,fill=BLUE] (p) circle[radius=0.1] node[anchor=south]{$p$};

	% R^mの座標系
	\draw[ORANGE!40, very thin, step=1] (15.1,6.5) grid(21.9,10.5);
	\draw[->,thick,color=ORANGE] (15,8) --(22,8) node[right]{$x_1$};
	\draw[->,thick,color=ORANGE] (18,6.5) --(18,10.5) node[right]{$x_2$};

	% Udash
	\begin{scope}[local bounding box=Udash]
		\draw[densely dashed,line width=1pt,color=ORANGE] (18.5, 8.5) circle[radius=1.5] (20, 10) node{$\varphi(U)$};
	\end{scope}

	% 2つの円の共通部分
	\begin{scope}
		\clip (18.5,8.5) circle[radius=1.5];
		\draw[densely dashed,line width=1pt,color=GREEN] (19.3,6.9) circle[radius=1.5];
	\end{scope}

	% 点phi(p)
	\coordinate (p) at (18.8, 7.7);
	\draw[color=BLUE,fill=BLUE] (p) circle[radius=0.1] node[anchor=west]{$\varphi(p)$};

	% R^mの座標系
	\draw[GREEN!40, very thin, step=1] (15.1,0.5) grid(21.9,4.5);
	\draw[->,thick,color=GREEN] (15,2) --(22,2) node[right]{$y_1$};
	\draw[->,thick,color=GREEN] (18,0.5) --(18,4.5) node[right]{$y_2$};

	% Vdash
	\begin{scope}[local bounding box=Vdash]
		\draw[densely dashed,line width=1pt,color=GREEN] (18.5, 2.5) circle[radius=1.5] (20, 2) node{$\psi(V)$};
	\end{scope}

	% 2つの円の共通部分
	\begin{scope}
		\clip (18.5,2.5) circle[radius=1.5];
		\draw[densely dashed,line width=1pt,color=ORANGE] (17.9,4.3) circle[radius=1.5];
	\end{scope}

	% 点psi(p)
	\coordinate (p) at (18.3, 3.4);
	\draw[color=BLUE,fill=BLUE] (p) circle[radius=0.1] node[anchor=east]{$\psi(p)$};

	% 繋ぎ線
	\draw[->,thick,shorten <=10pt, shorten >=10pt] (Udash.south) --(Vdash.north) node[midway, right]{$\psi\circ\varphi^{-1}$};
	\draw[->,thick,shorten <=10pt, shorten >=10pt] (Udash.west) --(U.east) node[midway, above]{$\varphi^{-1}$};
	\draw[->,thick,shorten <=10pt, shorten >=10pt] (Vdash.west) --(V.east) node[midway, below]{$\psi^{-1}$};

\end{tikzpicture}

	}
	\caption{座標変換の模式図}
	\label{fig:sample-coordinate-transform}
\end{figure}

\subsection{生成されるファイル}

\texttt{main.pdf}は, \texttt{main.tex}をコンパイルして得られるPDFである. 
また, コンパイル時には\texttt{aux/}や\texttt{main.log}, \texttt{main.synctex.gz}などの補助ファイルが生成されることがある. 
これらは本文やスタイルの本体ではなく, \LaTeX が処理のために作るファイルである. 

\end{document}
